\section{Method}\label{sec:method}

\paragraph{Model and fine-tuning} All experiments use AfroXLMR-base (\texttt{Davlan/afro-xlmr-base}), which is an XLM-R-base model adapted to African languages~\cite{alabi2022}, with a linear token-classification head. A single set of hyperparameters is used for every comparison: AdamW, a learning rate of $5\times10^{-5}$, weight decay of 0.01, batch size 16, 3 epochs, linear decay with no warm-up, gradient clipping at 1.0, fp16 and a maximum length of 128 subword tokens. Every configuration is trained with seeds 13, 42 and 2026 and the final model is used, therefore no checkpoint is selected on development or target data.

\paragraph{Zero-shot grid} Eleven sources, namely the six NCHLT languages, the three MasakhaNER target languages, and MasakhaNER Swahili and Chichewa, are trained with three seeds and evaluated on all nine test sets, which gives 297 zero-shot evaluations. The nine pairs that share language and corpus give in-language upper bounds, and the other 90 are transfer pairs. Each transfer pair is coded for \emph{relatedness} on four ordinal levels (same language in the other corpus, same cluster, other Southern Bantu and other Bantu, where higher means closer), \emph{domain match} (same corpus), \emph{target exposure} (seen or unseen), source--target \emph{entity overlap} and, in exploratory models, the \emph{fertility gap}, which is the absolute difference in fertility. Since each corpus has its own guideline, domain match bundles text domain with annotation conventions.

\paragraph{Few-shot transfer} For each news target (isiZulu, isiXhosa and Setswana), two zero-shot models are fine-tuned further on $k\in\{100,500\}$ target sentences, resampled with each seed: the sister-language news source (isiXhosa for isiZulu and vice versa) and the same-language government source. No Sotho-Tswana news data exists, so Setswana's sister source is NCHLT Sepedi. This gives 36 evaluations, or 333 in total, with 10 epochs at $k=100$ and 5 at $k=500$.

\paragraph{Measures} The primary measure is entity-level micro-F1 with exact span match on PER, LOC and ORG. The \emph{transfer ratio} is $\TR(s\!\rightarrow\!t)=F1(s\!\rightarrow\!t)/F1(t\!\rightarrow\!t)$, which normalises for target difficulty. For few-shot runs, the \emph{gap closure} at $k$ is $(F1_k-F1_0)/(F1_{\mathrm{in}}-F1_0)$ and the \emph{share of gain by 100} is $(F1_{100}-F1_0)/(F1_{500}-F1_0)$. All scores are means over three seeds, and the seed standard deviations are mostly under 3~F1 (up to 6 for NCHLT Sesotho).

\paragraph{Statistics} Five contrasts were pre-specified. For RQ1, the sister-language news source is compared with the same-language government source on isiZulu news and on isiXhosa news; for RQ4, Sepedi and Sesotho are compared with Swahili on Setswana news, and Sepedi with Swahili on Setswana government text. Each contrast is a paired bootstrap over test sentences with 1,000 resamples~\cite{dror2018}, where every resample draws the same sentences for all three seeds and compares seed-averaged F1. Tests are one-sided in the hypothesised direction with Holm correction over the five contrasts, and two-sided values are also reported. With 1,000 resamples the smallest attainable one-sided $p$ is .001 (two-sided .002), therefore the Holm values of .005 and .01 reported below are at that floor. Moreover, the bootstrap resamples test sentences and not training runs, so it does not capture the training variance between seeds.

A linear mixed-effects regression (M0) of transfer ratio on relatedness, domain match, their interaction, exposure and entity overlap is fitted on the 90 seed-averaged transfer pairs with a random intercept for the target language; the proposal also specified a random intercept for the source language, which the fitted model omits. RQ3 is answered descriptively, and Table~\ref{tab:rqmap} maps each question to its test. The robustness checks in Section~\ref{sec:res-robust} were added after the main results and are not part of the Holm family.

\begin{table}[t]
\centering\footnotesize
\caption{How each research question is answered.}
\label{tab:rqmap}
\begin{tabular}{@{}l>{\raggedright}p{3.7cm}>{\raggedright\arraybackslash}p{2.6cm}@{}}
\toprule
RQ & Comparison & Measure and test\\
\midrule
RQ1 & Relatedness and domain match over 90 transfer pairs; sister-language news vs same-language government source on news targets & F1, TR; 2 bootstrap contrasts; M0\\
RQ2 & Seen (zul, xho, sot) vs unseen (ssw, nso, tsn) targets & F1, TR; exposure term in M0\\
RQ3 & 2 sources $\times$ 3 news targets, $k=0,100,500$ & F1, gap closure\\
RQ4 & Sepedi, Sesotho vs Swahili, Chichewa on both Setswana tests & F1; 3 bootstrap contrasts\\
\bottomrule
\end{tabular}
\end{table}
